\documentclass{article}
\usepackage{amsmath}
\usepackage{graphicx}

\title{Sample LaTeX Document}
\author{Test Author}
\date{\today}

\begin{document}

\maketitle

\section{Introduction}

This is a test LaTeX document to demonstrate the reader tool functionality.

\subsection{Motivation}

Lorem ipsum dolor sit amet, consectetur adipiscing elit. Sed do eiusmod tempor incididunt ut labore et dolore magna aliqua.

\section{Mathematical Background}

Consider the following equation:

\begin{equation}
E = mc^2
\end{equation}

This is Einstein's famous mass-energy equivalence formula.

\subsection{More Complex Math}

We can also write systems of equations:

\begin{align}
x + y &= 10 \\
2x - y &= 5
\end{align}

\section{Experimental Setup}

Our experimental setup consists of several components.

\begin{figure}
\centering
\includegraphics[width=0.8\textwidth]{example-image}
\caption{Experimental setup diagram}
\label{fig:setup}
\end{figure}

As shown in Figure \ref{fig:setup}, the system architecture is straightforward.

\section{Results}

The results of our experiments are summarized in Table \ref{tab:results}.

\begin{table}
\centering
\begin{tabular}{|l|c|r|}
\hline
Method & Accuracy & Time (s) \\
\hline
Baseline & 85\% & 10.5 \\
Proposed & 95\% & 8.2 \\
\hline
\end{tabular}
\caption{Performance comparison}
\label{tab:results}
\end{table}

\subsection{Analysis}

The improved performance can be attributed to our novel approach.

\section{Conclusion}

In this paper, we have demonstrated the effectiveness of our method.

\subsection{Future Work}

Future research directions include:
\begin{itemize}
\item Extending to larger datasets
\item Optimizing runtime performance
\item Exploring alternative algorithms
\end{itemize}

\end{document}
